\section{Exact arithmetic conventions}\label{app:arithmetic}
All constants in the certificates are rational; a terminating decimal denotes the corresponding rational number exactly.
This appendix explains the common sign-certification rules.
The individual polynomials, parameter maps, and subdivision orders are specified where they are used.

\subsection{Bernstein coefficients}
We use the standard enclosure supplied by the Bernstein form of a
polynomial; see Cargo and Shisha~\cite[Section~3]{CargoShisha66} for
the univariate coefficient identity and range bound. We include the
tensor-product formula and its proof to make the rational checks explicit.
Let \(P(x,y)=\sum_{k=0}^N\sum_{\ell=0}^K a_{k\ell}x^ky^\ell\).
On the unit square its Bernstein expansion is
\begin{equation}\label{eq:bernstein-expansion}
P(x,y)=\sum_{i=0}^N\sum_{j=0}^K b_{ij}
\binom Ni x^i(1-x)^{N-i}\binom Kj y^j(1-y)^{K-j},
\end{equation}
where
\begin{equation}\label{eq:bernstein-coefficients}
b_{ij}=
\sum_{\substack{0\le k\le i\\0\le\ell\le j}}
a_{k\ell}
\frac{\binom ik}{\binom Nk}
\frac{\binom j\ell}{\binom K\ell}.
\end{equation}
Indeed,
\[
x^k=\sum_{i=k}^N
\frac{\binom ik}{\binom Nk}\binom Ni x^i(1-x)^{N-i},
\]
by the binomial theorem after factoring out \(x^k\); multiplying two such identities proves the formula.
The basis functions in Equation~\eqref{eq:bernstein-expansion} are nonnegative and sum to one.
Therefore
\begin{equation}\label{eq:bernstein-lower}
P(x,y)\ge\min_{i,j}b_{ij}\qquad(0\le x,y\le1).
\end{equation}
The univariate version is obtained by omitting \(y\).

For a rational rectangle, first substitute the affine maps from the unit square.
For a polynomial obtained after a rational or algebraic change of variables, the displayed denominator is cleared only where it is positive, as verified in the corresponding argument.
A positive common scale may be applied to all coefficients, making every calculation integral.
If a table reports
\(\lfloor C\min b_{ij}\rfloor\) for a stated \(C>0\), a positive table entry is an exact strict sign certificate.

Subintervals can be handled either by direct substitution or by de Casteljau subdivision.
For subdivision at \(1/2\), start from a Bernstein row
\((b_0,\ldots,b_N)\), repeatedly replace neighboring entries by their arithmetic mean, and take the left and right outer diagonals as the Bernstein rows on the two halves.
This rule follows by substituting \(x=t/2\) and \(x=(1+t)/2\) into the basis expansion.
It is exact over the rationals; a common power of two clears all new denominators.
For rectangles, apply the rule separately in each coordinate.

\subsection{Intervals and derivative enclosures}
An interval \([a,b]\) encloses a real quantity if \(a\le x\le b\).
Addition, subtraction, multiplication, and reciprocal use the endpoint extrema, with a reciprocal taken only when zero is excluded.
Square roots are taken only on nonnegative intervals.
For a positive integer scale \(Q\), rounding a lower endpoint down to a multiple of \(1/Q\) and an upper endpoint up preserves enclosure.
Scaled square roots reduce to integer square-root bounds:
for an integer \(A\ge0\),
\[
\lfloor\sqrt A\rfloor^2\le A<
(\lfloor\sqrt A\rfloor+1)^2.
\]
Consequently no floating-point operation is needed for an interval certificate.

A derivative enclosure supplies a second bound on a box.
If a function has center enclosure \(I_0\), coordinate half-widths \(h_i\), and derivative magnitude bounds \(M_i\), then
\begin{equation}\label{eq:mean-value-box}
P(\text{box})\subseteq
I_0+\left[-\sum_i h_iM_i,\;\sum_i h_iM_i\right].
\end{equation}
This follows by integrating the derivative along a segment from the center.
The same statement holds with bounded limiting derivatives at a boundary whenever the resulting Lipschitz estimate extends continuously.
At a square-root singularity the direct interval enclosure is used instead.
The pair certificate in Appendix~\ref{app:scalar} specifies which bounds are used and how uncertified boxes are subdivided.
An undefined local interval operation causes further subdivision, never acceptance.

\subsection{Supplementary sources}
The complete certificate sources accompany the paper in the directory
\texttt{verification/}.
They are organized as follows.
\begin{center}
\small
\begin{tabular}{@{}lp{.52\textwidth}@{}}
\toprule
Source & Certified calculations\\
\midrule
\texttt{interval\_certificate.cpp}
 & The three pair cases in Appendix~\ref{app:scalar}.\\
\texttt{profile\_certificate.py}
 & The polynomial profiles, their branch inequalities, and the three-bit base cases.\\
\texttt{central\_band.py}
 & The central-band polynomial signs and rational constant bounds.\\
\texttt{intermediate\_bias.py}
 & The six intermediate-bias Bernstein bounds.\\
\texttt{large\_bias.py}
 & The four norm-estimate bounds, including the series tail.\\
\texttt{profile\_formula\_audit.py}
 & An independent coefficient reconstruction of the profile formulas and printed data.\\
\bottomrule
\end{tabular}
\end{center}
The accompanying README gives the compiler and Python requirements and a single command to reproduce the calculations.
The C++ source uses outward-rounded integer intervals with signed extended 128-bit intermediates; Appendix~\ref{app:scalar} gives the relevant bounds.
The principal Python profile certificate uses arbitrary-precision integer arithmetic.
The symbolic reconstruction scripts use exact rational coefficients.
The arguments proving the enclosure rules, the complete domain coverage, and the passage from the scalar inequalities to the dimension-free theorem are part of the paper.

